\chapter{Conclusion and Future Work}

\begin{figure}[h]
\centering
\includegraphics[width=\linewidth]{../figs/schematic.png}
\caption{Summary of changes in physical drivers and outcomes of the 2025 Sumatra extreme rainfall. Values indicate the relative change for the historical attribution (past-to-present) and future projection (present-to-future) scenarios. Blue boxes highlight localized impacts within the Aceh Province, while gray boxes show regional changes across the Malacca Strait, Sumatra, and Malay Peninsula. Heavy rainfall intensity and high updraft is defined by the 99.9th percentile of the variables in factual simulation, and the rainfall area represents the geographic extent exceeding the 300 mm threshold over the peak two-day period. Similary, latent heating area and high-IVT area represent the geographic extend exceeding its 99th percentile threshold in the factual simulation over the peak period. Evaporation and CAPE changes are from the distribution median change.}
\label{fig:schematic}.
\end{figure}

This study employed a high-resolution, convection-permitting RegCM5 model and the pseudo-global warming (PGW) framework to attribute the influence of anthropogenic climate change (ACC) to the rare equatorial Tropical Cyclone (TC) Senyar and the resulting 2025 Sumatra flash floods. Our reconstruction of the event confirms that the high-resolution model provides significant value over coarse reanalysis data (ERA5), particularly in capturing the spatial distribution of extreme accumulated rainfall across northern Sumatra. While the presented attribution focuses on Aceh Province, similar findings are found in the other affected provinces, such as we argue that the findings are applicable for Sumatra in general.

The attribution analysis reveals a distinct change in the hazard's character across three different climate (past, present, and future):
\begin{itemize}
    \item Historical warming has intensified the 2025 event. While the total rainfall in Aceh increased by 10\%, the most striking finding is a 64\% expansion in the rainfall area exceeding 300 mm. This indicates that climate change has drastically increased the geographic exposure to extreme hydrological triggers compared to a pre-industrial climate.
    \item In a +1.5$^\circ$ warmer world, the hazard may become more spatially concentrated but significantly more violent. While heavy rainfall intensity continues to increase (+17\%), exceeding Clausius-Clapeyron scaling, the total rainfall and total rainfall area are projected to contract by 9\% and 15\%, respectively.
    \item These changes are underpinned by a robust increase in atmospheric energy and moisture transport. We identify more than 20\% increase in CAPE across scenarios and an intensification of extreme IVT (+5\%–12\% per degree) specifically within the moisture corridor of the Malacca Strait. As shown in the summary schematic (Figure \ref{fig:schematic}), the release of increased latent heat within the storm core accelerates high updraft velocities (more than 10\% of increase), effectively converting moisture to concentrated bursts even as the broader storm footprint shrinks.
\end{itemize}

Finally, while TC Senyar’s trajectory remains generally robust across all experiments, the projection of increased peak wind intensity in the future scenario suggests that Sumatra may face an emerging threat of compounding wind and hydrological hazards as the climate continues to warm.

To improve the attribution of the 2025 floods, several aspects for future research remain:
\begin{itemize}
    \item We can project the extreme event in a warmer climate, e.g +3.0$^\circ$ global warming level. This will complete our understanding of the non-linear hazard responses under higher warming scenarios.
    \item A primary limitation of the current study was the inability to perfectly represent the daily rainfall evolution at individual meteorological stations, partly due to UTC-to-local time offsets and the inherent challenges of simulating localized convective timing. Future work should prioritize verifying these results with high-resolution hourly observation data once it becomes available.
    \item The 2025 disaster was a result of both atmospheric triggers and terrestrial vulnerability, especially with the known drastic deforestation and the unique geographical feature of hills in Sumatra \citep{wu2026analysis, wwa}. Future studies should integrate these amplified rainfall projections with land-use change to quantify how deforestation, floodings, and landslide susceptibility interact with the rainfall and wind intensity response.
    \item Expanding this storyline approach to a larger ensemble of CMIP6 models would help further quantify the structural uncertainties in the evolution of TC Senyar and extreme rainfall and provide more constrained risk assessments for local government’s adaptation strategies.
\end{itemize}